% Adición acumulativa. Insertar después de la suspensión esturmiana existente.
% Conservar íntegro el texto y la figura fig:holonomia-esturmiana.
\subsection{Conexión nonádica, holonomía y representación monodrómica}
\label{mono-ampl:conexion}

La conexión regional ya construida se estudia ahora mediante tres objetos:
el transporte del estado completo, la representación de sus vueltas y la
suspensión de sus itinerarios. Las operaciones de APP--TRIT--TPK y el
catálogo inicial conservan las definiciones anteriores. La distinción
necesaria concierne aquí a lo que cada representación retiene de una
misma historia.

En el índice de prolongación \(k\), una coordenatización del estado conserva
\[
 \widehat x_k=(B_k,C_k,p_k,c_k,\Sigma_k,W_k,g_k,m_k).
\]
Aquí \(B_k\) reúne los tres bloques regionales, \(C_k\) es el cilindro,
\(p_k\) el prefijo, \(c_k\) el registro de acarreo, \(\Sigma_k\) la firma
de frontera y \(W_k\) la información incidencial. La fase es
\(g_k=1+[(k-1)]_9\); \(m_k\) cuenta las vueltas completas. La reducción
\(\beta_k=[m_k]_{12}\) conserva la posición dodecafásica, mientras el
entero \(m_k\) distingue sus sucesivas vueltas.

\begin{definition}[Transporte de una vuelta nonádica]
\label{mono-ampl:def-vuelta}
Sean \(\mathcal R_k\subseteq\widehat X_k\times\widehat X_{k+1}\)
las relaciones de extensión admisible del TPK. La realización de la
holonomía en el nivel \(k\) es la composición relacional
\[
 \Gamma_{9,k}=\mathcal R_{k+8}\circ\cdots\circ\mathcal R_k.
\]
Para \(s\geq1\), su prolongación a \(s\) vueltas es
\[
 \Gamma_{9s,k}=
 \Gamma_{9,k+9(s-1)}\circ\cdots\circ\Gamma_{9,k}.
\]
\end{definition}

La formulación relacional conserva las continuaciones admisibles. Sobre
una historia individualizada, cada flecha lleva al estado efectivamente
prolongado; el bloque visible aislado puede seguir admitiendo más de una
continuación. Esta diferencia es precisamente la función de la memoria y
de las condiciones de frontera en la determinación de la trayectoria.

\begin{proposition}[Iteración del retorno y memoria acumulada]
\label{mono-ampl:vueltas}
En toda composición admisible de \(s\) vueltas,
\[
 g_{k+9s}=g_k,\qquad m_{k+9s}=m_k+s,\qquad
 \beta_{k+9s}=\beta_k+s\pmod{12}.
\]
\end{proposition}
\begin{proof}
La regla de una vuelta conserva la fase e incrementa el contador en una
unidad. Su aplicación al estado ya transportado repite exactamente esas
dos operaciones. La inducción da las dos primeras igualdades; reducir la
segunda módulo \(12\) da la tercera.
\end{proof}

El contador proporciona así un testigo entero de la diferencia entre las
vueltas. Tras \(12\) vueltas la pareja de fases finitas se restaura, pero
el contador aumenta \(12\). La extensión no escindida
\eqref{eq:extension108} expresa la composición del calendario mediante su
cociclo; la trayectoria conserva además la elevación entera de ese dato.
La conservación significa aquí transporte de las contribuciones previas,
no inmovilidad de todas las coordenadas.

\subsubsection{Representación bilateral del índice de retorno}

El desplazamiento \(T\) de \eqref{eq:weyl-relojes} admite una
expresión en coordenadas de fase y vuelta. Numeramos las fases de \(0\) a
\(8\), mediante un desplazamiento del origen respecto de \(g_k\), y
escribimos \(k=9m+r\), con \(0\leq r<9\), en la recta entera
que recubre el ciclo de fases. Sobre
\(\mathcal H_{\rm ind}=\ell^2(\mathbb Z)\), la base
\(|k\rangle=|r,m\rangle\) expresa ese mismo desplazamiento unitario
\[
 T|r,m\rangle=
 \begin{cases}
 |r+1,m\rangle,&r<8,\\
 |0,m+1\rangle,&r=8.
 \end{cases}
\]
La inversa reduce \(r\) en una unidad cuando \(r>0\), y lleva
\(|0,m\rangle\) a \(|8,m-1\rangle\). El paso por el origen de fase
incorpora, por tanto, la unidad exacta de cociente.

\Needspace{8\baselineskip}
\begin{proposition}[Representación de la monodromía del ciclo]
\label{mono-ampl:representacion}
Sea \(S=T^9\). Entonces
\[
 S|r,m\rangle=|r,m+1\rangle,\qquad
 \mathbb Z\longrightarrow\mathcal U(\mathcal H_{\rm ind}),
 \quad s\longmapsto S^s
\]
es una representación unitaria fiel del grupo de vueltas del ciclo.
\end{proposition}
\begin{proof}
Nueve desplazamientos conservan el resto de la división por \(9\) e
incrementan el cociente. Las potencias satisfacen
\(S^{s+t}=S^sS^t\), y \(S^s|r,m\rangle=|r,m+s\rangle\) distingue
todo entero \(s\). La permutación de una base ortonormal es unitaria.
\end{proof}

Esta representación corresponde al recubrimiento del ciclo, cuyo grupo
de vueltas es \(\mathbb Z\), y representa la fase con su contador. El
estado del TPK contiene además las coordenadas de trayectoria ya descritas.
La extensión bilateral admite índices negativos como coordenadas del
recubrimiento; una trayectoria iniciada en una condición inicial utiliza
su semieje admisible. De este modo se dispone de una inversa operatoria
del índice sin convertir cada relación de extensión en una biyección del
estado completo.

Con \(D_M|k\rangle=\lfloor k/9\rfloor|k\rangle\), la identidad
\([D_M,S]=S\) expresa el incremento de memoria en el dominio denso de
vectores de soporte finito. Si \(\mathcal I|k\rangle=|-k\rangle\) y
\(P_0\) proyecta sobre los índices divisibles por \(9\), se obtiene
además
\begin{equation}
 \mathcal I T\mathcal I=T^{-1},\qquad
 \mathcal I D_M\mathcal I=-D_M-(I-P_0).
 \label{mono-ampl:inversion-contador}
\end{equation}
En efecto, para \(k=9m+r\),
\(\lfloor-k/9\rfloor=-m\) si \(r=0\) y vale \(-m-1\) si \(r>0\).
La inversión conserva así el término de frontera de la división euclídea.

\subsection{Reloj de bloques y reloj de capacidad}
\label{mono-ampl:relojes}

El índice de una aplicación de refinamiento y el número de posiciones
decimales disponibles tienen leyes relacionadas, pero diferentes. Para
expresarlas reservamos \(n\) al índice de bloques ternarios y definimos
\[
 \rho=\log_{1000}729=\log_{10}9,\qquad
 \delta=1-\rho,\qquad
 \mathcal C_n=\lfloor(n+1)\rho\rfloor,\quad n\geq0.
\]
La letra \(\mathcal C_n\) designa aquí la capacidad, no el registro
dodecafásico \(K\). En el origen de fase canónico, la palabra mecánica
anterior proporciona
\[
 Z_n=\sum_{j=0}^{n}z_j(0)=\lfloor(n+1)\delta\rfloor.
\]

\begin{proposition}[Balance exacto entre capacidad y vacancias]
\label{mono-ampl:balance-capacidad}
Para \(n\geq0\) y, en la segunda identidad, \(n\geq1\),
\[
 \mathcal C_n=n-Z_n,\qquad
 \mathcal C_n-\mathcal C_{n-1}=1-z_n(0).
\]
\end{proposition}
\begin{proof}
La irracionalidad de \(\delta\) implica
\(\lceil(n+1)\delta\rceil=\lfloor(n+1)\delta\rfloor+1\).
Entonces
\[
 \lfloor(n+1)(1-\delta)\rfloor
 =n+1-\lceil(n+1)\delta\rceil=n-Z_n.
\]
Restar dos identidades consecutivas concluye la prueba.
\end{proof}

Un evento \(z_n=1\) añade un bloque a la historia manteniendo la capacidad
\(\mathcal C_n\); un evento \(z_n=0\) incrementa ambos contadores. En
cualquier intervalo, los incrementos de capacidad y las vacancias suman
exactamente su longitud. Reduciendo módulo \(q\), para \(q=9\) o
\(q=108\), resulta
\[
 [\mathcal C_n]_q=[n]_q-[Z_n]_q.
\]
La fase uniforme de bloques y la fase de capacidad se relacionan mediante
la memoria integrada de vacancias. En particular, una retención de
capacidad y una vuelta completa de la conexión tienen mecanismos
distintos, aunque ambas puedan conservar una observación de fase.

La inversión monótona del reloj da un tiempo de primera llegada:
\[
 T_{\rm cap}(a)=\min\{n:\mathcal C_n=a\}
 =\left\lceil\frac a\rho\right\rceil-1,
 \qquad a\geq1.
\]
Con \(\sigma=\rho^{-1}-1\), el tiempo empleado en ganar \(b\) unidades
de capacidad es
\[
 T_{\rm cap}(a+b)-T_{\rm cap}(a)
 =b+\lceil(a+b)\sigma\rceil-\lceil a\sigma\rceil.
\]
La diferencia de techos toma uno de los dos enteros adyacentes a
\(b\sigma\). Así se obtienen \(9\) o \(10\) bloques para ganar
\(9\) unidades de capacidad, y \(113\) o \(114\) para ganar \(108\).
Son tiempos del reloj inverso, medidos desde una primera llegada; el
contador de vueltas se interpreta siempre en el índice al que pertenece.

\subsection{Suspensión simbólica y cristal temporal aperiódico}
\label{mono-ampl:cristal}

La codificación por rotación reúne periodicidad de fase, recurrencia local
y aperiodicidad global. Sea \(\mathcal X_\delta\subseteq\{0,1\}^{\mathbb Z}\)
la clausura, en la topología producto, de las traslaciones de la palabra
bilateral \(z_n(\theta)\). Denotemos por \(\sigma_{\rm sh}\) el
desplazamiento de sus sucesiones. Conservando una fase nonádica, el paso es
\[
 F(r,z)=([r+1]_9,\sigma_{\rm sh}z),\qquad
 (r,z)\in C_9\times\mathcal X_\delta.
\]

\begin{definition}[Cristal temporal aperiódico simbólico]
\label{mono-ampl:def-cristal}
La suspensión de techo unitario del sistema anterior es el espacio
\[
 \mathcal S_\delta=
 \bigl((C_9\times\mathcal X_\delta)\times\mathbb R\bigr)/\sim,
 \qquad (y,t+1)\sim(Fy,t),
\]
con flujo \(\Phi_s[y,t]=[y,t+s]\). Se denomina \emph{cristal temporal
aperiódico simbólico} a este modelo de orden temporal, junto con su
codificación de vacancias y la representación del contador de vueltas.
\end{definition}

El adjetivo temporal designa la acción del flujo y de su sección de retorno;
la aperiodicidad caracteriza sus itinerarios. La frecuencia de unos de cada
sucesión es \(\delta\): la cota uniforme de discrepancia menor que uno,
obtenida por telescopía en cualquier ventana, persiste al tomar la clausura.
Una sucesión periódica tendría frecuencia racional. Por tanto,
\(\mathcal X_\delta\) no contiene órbitas periódicas. La suspensión
tampoco tiene órbitas cerradas: un retorno del flujo exigiría un tiempo
entero y una periodicidad del itinerario.

Las palabras finitas sí reaparecen. Los itinerarios de longitud \(m\)
proceden de una partición de la circunferencia por \(m+1\) puntos; cada
intervalo es visitado recurrentemente por la rotación irracional. Incluso
al restringir los tiempos a una fase fija, el paso \(9\delta\) sigue
siendo irracional. Ésta es la razón de que cada palabra aparezca en las
\(9\) fases y de las complejidades ya establecidas
\[
 p(m)=m+1,\qquad p_9(m)=9(m+1),\qquad p_3(m)=3(m+1).
\]
La recurrencia de configuraciones locales coexiste con la ausencia de un
período global. Su entropía topológica es nula porque estas complejidades
crecen linealmente.

La realización por los observables \(V_9,W_\delta\) y el desplazamiento
\(T\) expresa el mismo orden mediante las covariancias
\eqref{eq:weyl-relojes}. Las frecuencias del reloj uniforme pertenecen a
\(\frac19\mathbb Z+\delta\mathbb Z\), módulo \(1\). El reloj de
capacidad conserva adicionalmente la relación
\([\mathcal C_n]_9=[n]_9-[Z_n]_9\); no se sustituye por la fase uniforme
al interpretar los retornos. El término \emph{cristal temporal} queda así
definido mediante un sistema dinámico matemático. Su realización física
requiere especificar una dinámica energética, sus observables temporales
y su estabilidad, datos distintos de la construcción simbólica presente.

\subsection{Simetría especular de las regiones y doble realización}
\label{mono-ampl:regiones}

El intercambio especular actúa en el bloque regional
\(B=(u^\pi,u^e,u^\varphi)\in(\mathbb F_3^6)^3\) por
\[
 \mathfrak cB=(u^e,u^\pi,u^\varphi).
\]
Su eje fijo y su agregado son, respectivamente, \(u^\varphi\) y
\(s=u^\pi+u^e\). Esta fijación se refiere a la involución en una fibra;
ambas coordenadas pueden evolucionar durante el transporte nonádico.
Con \(q=u^\pi+u^e+u^\varphi\),
\(a=u^\pi+u^e-u^\varphi\) y \(d=u^\pi-u^e\), se recupera
\[
 u^\varphi=2(q-a),\qquad s=q-u^\varphi,\qquad
 u^\pi=2(s+d),\qquad u^e=2(s-d).
\]
Las igualdades son componente a componente en \(\mathbb F_3\), donde
\(2^{-1}=2\). El espejo conserva \(q,a,s\) e invierte \(d\); esta
coordenada antisimétrica individualiza el reparto que el agregado conserva
sin distinguir. La inversión \(\mathcal I\) del índice temporal y la
involución \(\mathfrak c\) de canales actúan, por consiguiente, sobre
datos diferentes.

El retorno de fase posee un testigo explícito:
\[
 \begin{aligned}
 B_1&=(012222\mid121221\mid112202),\\
 B_{10}&=(101222\mid112221\mid112002).
 \end{aligned}
\]
Las dos posiciones tienen fase \(1\); el bloque, el eje y la memoria han
cambiado. Asimismo, la sincronización de los censos en las fases \(8\) y
\(9\) anula sus diferencias relativas, mientras la componente diagonal
pasa de \((59,59,59)\) a \((43,43,43)\). La igualdad de un observable
relativo no agota el estado que lo determina.

La doble realización conserva esa arquitectura. En la sección
\ref{exc:doble-lectura}, el mismo prefijo determina por una parte cilindros
encajados y por otra una sucesión de palabras codificadas con sus marcos.
El primer mapa pasa al límite por contracción de diámetros; el segundo
satisface \(r_{n,m}Q_n=Q_mp_{n,m}\) y define \(Q_\infty\).
La figura \ref{mono-ampl:fig-regional} anticipa esas dos aplicaciones.
El recorrido incidencial posterior construye códigos, retículos y la
realización de Moonshine en sus dominios propios; el grupo de automorfismos
actúa sobre esa realización, no sobre las cifras por una identificación
implícita.

% Figura acumulativa; destino figures/monodromia_regional_ampliacion.tex.
\begin{figure}[!htbp]
\centering
\begin{tikzpicture}[x=1cm,y=1cm,>=Latex,font=\small,
  flow/.style={->,line width=.55pt},
  ann/.style={align=center,font=\footnotesize}]
\node[font=\bfseries] at (6.9,9.2)
  {Conexión nonádica y coordenadas regionales};
\begin{scope}[shift={(3.4,6.8)}]
 \foreach \k in {1,...,9}{
  \pgfmathsetmacro{\ang}{90-40*(\k-1)}
  \coordinate (p\k) at (\ang:1.52);
  \fill (p\k) circle (1.4pt);
  \node[fill=white,inner sep=1.1pt,font=\footnotesize]
    at (\ang:1.83) {\(\k\)};
 }
 \foreach \k in {1,...,8}{
  \pgfmathtruncatemacro{\next}{\k+1}
  \draw[flow,shorten >=3pt,shorten <=3pt] (p\k)--(p\next);
 }
 \draw[flow,shorten >=3pt,shorten <=3pt] (p9)--(p1);
 \node[ann] at (0,.1) {\(g_k\in C_9\)\\\(\Gamma_{9,k}\)};
 \node[ann] at (0,-2.23) {Una vuelta: \(g\mapsto g\), \(m\mapsto m+1\)};
\end{scope}
\node[ann,text width=5.25cm] at (10.1,7.75)
 {\(B_k=(u_k^\pi,u_k^e,u_k^\varphi)\)\\
 bloque regional en una fibra};
\node at (8.55,6.5) {\(u^\pi\)};
\node at (11.65,6.5) {\(u^e\)};
\draw[<->,line width=.65pt] (8.95,6.5)--node[above]{\(\mathfrak c\)}(11.25,6.5);
\draw[densely dashed,gray] (10.1,5.72)--(10.1,7.23);
\node[fill=white,inner sep=2pt] at (10.1,5.6){\(u^\varphi\)};
\node[ann] at (10.1,4.6)
 {Eje \(u^\varphi\) y agregado \(u^\pi+u^e\)\\
 invariantes bajo \(\mathfrak c\)};
\draw[flow] (5.45,6.8)--(7.25,6.8);
\draw[gray!60,line width=.3pt] (.55,4.02)--(13.25,4.02);
\node[ann,text width=12cm] (hist) at (6.9,3.47)
 {Historia compatible: prefijos, cilindros, orientación, cocientes,
 marcos de incidencia y memoria};
\coordinate (fork) at (6.9,2.85);
\draw[flow] (hist.south)--(fork);
\draw[flow] (fork)--(3.6,2.2);
\draw[flow] (fork)--(10.3,2.2);
\node[ann,text width=5.6cm] at (3.6,1.65)
 {Realización arquimediana\\
 \(I_{n+1}\subseteq I_n,\quad |I_n|=729^{-n}\)\\
 límite de los prefijos regionales};
\node[ann,text width=5.6cm] at (10.3,1.65)
 {Realización incidencial\\
 \(r_{n,m}Q_n=Q_mp_{n,m}\)\\
 palabras codificadas y marcos compatibles};
\node[ann,text width=5.6cm] at (3.6,.42)
 {\(\pi,\varphi,e\); determinación dodecafásica de \(\alpha\)};
\node[ann,text width=5.6cm] at (10.3,.42)
 {Witt--Golay--Leech--\(V^\natural\)\\
 automorfismos de la realización excepcional};
\draw[flow] (3.6,1.03)--(3.6,.78);
\draw[flow] (10.3,1.03)--(10.3,.78);
\end{tikzpicture}
\caption{\fontsize{9}{11}\selectfont Nueve fases de una conexión y dos realizaciones de la historia
compatible. El ciclo representa la fase; cada vuelta incrementa el contador.
El diagrama regional representa la involución \(\mathfrak c\), que
intercambia clausura y propagación y fija el eje de autoescala en una fibra.
Esta invariancia no impone que el eje sea constante durante el transporte.
Las dos aplicaciones inferiores utilizan el mismo prefijo: una determina
su límite numérico y la otra conserva palabras y marcos de incidencia. Los
mapas de esta segunda realización se especifican en la sección
\ref{exc:doble-lectura}; el recorrido excepcional se desarrolla después.}
\label{mono-ampl:fig-regional}
\end{figure}


\clearpage
\subsubsection{Representación rectilínea de las fases y del retorno}
\label{mono-recta:desarrollo}

La representación sobre el segmento \([0,10]\) conserva el origen
posicional de APP y ordena sus nueve marcas interiores. En este diagrama,
los enteros \(1,\ldots,9\) son representantes de fase. La posición
horizontal indica el orden de las coordenatizaciones; las etiquetas superiores
especifican operaciones regionales. La figura
\ref{mono-recta:figura} reúne así la recta de partida, la fijación del
eje de autoescala, la distribución especular de los dos canales y el
retorno con memoria.

% Recuperación de fig_nonadic_route_mini.tex de la síntesis académica.
% Los rótulos designan operaciones regionales; la suma se efectúa en F_3^6.
\begin{figure}[H]
\centering
\begin{tikzpicture}[x=1.12cm,y=1cm,font=\small,>=Latex,
  guide/.style={line width=.45pt,black!55}]
 \draw[line width=.8pt] (0,0)--(10,0);
 \foreach \k in {0,...,10}{
  \draw[guide] (\k,-.09)--(\k,.09);
  \node[below=2.5mm] at (\k,0) {$\k$};
 }
 \foreach \k in {1,...,9}{\fill (\k,0) circle (1.7pt);}
 \node[anchor=west,font=\footnotesize] at (0,3.05)
   {Segmento generador $[0,10]$ y nueve representantes de fase};
 \draw[guide] (5,.12)--(5,1.7);
 \node[above,align=center,font=\footnotesize] at (5,1.7)
   {fijación del eje\\$u^\varphi$};
 \draw[decorate,decoration={brace,amplitude=4pt},guide]
   (6,.65)--(8,.65);
 \node[above=5pt,align=center,font=\footnotesize] at (7,.65)
   {agregado regional\\$u^e+u^\pi$};
 \draw[guide] (9,.12)--(9,1.7);
 \node[above,align=center,font=\footnotesize] at (9,1.7)
   {involución especular\\$\pi\leftrightarrow e$};
 \draw[->,line width=.7pt] (9,-.7)
   .. controls (9,-1.55) and (1,-1.55) .. (1,-.7);
 \node[fill=white,inner sep=3pt,font=\footnotesize] at (5,-1.23)
   {retorno $9\to1$;\quad $m\mapsto m+1$};
 \node[font=\footnotesize] at (5,-1.95)
   {$w_6\to w_{12}\to w_{18}\to w_{24}\to w_{30}\to R_{36}\to G_9$};
\end{tikzpicture}
\caption{Representación rectilínea de la conexión nonádica. La marca
\(5\) sitúa la primera saturación fuerte y el eje regional de
autoescala; la llave \(6\)--\(8\) identifica el agregado ternario
fijado en cada fibra; la fase \(9\) selecciona la orientación de
clausura. El retorno conserva la fase e incrementa la memoria. Las
abscisas numeran posiciones operativas, mientras las etiquetas
\(u^\pi,u^e,u^\varphi\) designan bloques de \(\mathbb F_3^6\).}
\label{mono-recta:figura}
\end{figure}


La fijación del eje se formula en la fibra de una frontera dada. Si
\(q=u^\pi+u^e+u^\varphi\) y
\(a=u^\pi+u^e-u^\varphi\), entonces
\[
 u^\varphi=2(q-a),\qquad s=u^\pi+u^e=q-u^\varphi
 \quad\text{en }\mathbb F_3^6.
\]
Los datos \((q,a)\) determinan el eje y el agregado; el reparto ordenado
de \(s\) conserva la libertad especular resuelta por la prolongación.
La quinta fase posee una solución local y una extensión, con dato
de frontera \(c=111111\) y firma residual \(\rho_W=000\). Éste es el
sentido de su primera saturación fuerte. El símbolo \(\varphi\)
situado sobre la marca \(5\) identifica el eje regional que interviene
en esa operación.

Las fases intermedias conservan esta descomposición dentro de cada fibra,
mientras sus datos de frontera evolucionan. El censo completo de
multiplicidades locales y extensiones de la rama distinguida es
\[
\begin{array}{c|rrrrrrrrr}
\text{fase}&1&2&3&4&5&6&7&8&9\\\hline
N_{\rm loc}&3&2&5&4&1&1&3&3&2\\
N_{\rm ext}&2&5&4&1&1&3&3&2&1
\end{array}
\]
La fase \(6\) combina unicidad local y tres prolongaciones; la fase
\(7\) resuelve una fibra especular triple; la fase \(8\) sincroniza
los tres censos regionales. La fase \(9\) conserva dos distribuciones
locales con el mismo eje y el mismo agregado, de las cuales una se
prolonga. Estos recuentos expresan funciones diferentes dentro de una
misma conexión, y justifican las marcas de la recta.

En particular, para las fases \(6,7,8\) las sumas regionales son
respectivamente \(012100\), \(101001\) y \(112000\). La conservación
representada por la llave es una invariancia respecto del reparto
especular en cada fibra. La evaluación real posterior aplica los
lectores de los cilindros compatibles a las historias completas; la
suma interna representada en la figura pertenece al dominio ternario.

Finalmente, el arco \(9\to1\) representa el paso de una vuelta a la
siguiente. Las fórmulas de la sección
\ref{mono-ampl:conexion} conservan el resto de fase e incrementan el
contador de vueltas. El diagrama rectilíneo y el diagrama cíclico de la
figura \ref{mono-ampl:fig-regional} expresan, por tanto, dos aspectos
complementarios de la misma holonomía: orden de las posiciones y
composición del retorno sobre el estado con memoria.


\subsection{Transporte temporal y composición exponencial trítica}
\label{mono-ampl:euler}

Las relaciones de Euler de la sección \ref{euler-trit-articulacion}
describen las realizaciones locales de los tres regímenes. Para una firma
fija \(\tau\), el generador satisface \(J_\tau^2=-\tau I\); la
conjugación algebraica lleva \(J_\tau\) a \(-J_\tau\). En consecuencia,
\[
 \exp(sJ_\tau)\exp(tJ_\tau)=\exp((s+t)J_\tau),\qquad
 \overline{\exp(tJ_\tau)}=\exp(-tJ_\tau),
\]
y el producto de una evolución con su conjugada es la unidad. En el tipo
elíptico se realiza mediante \(\cos^2t+\sin^2t=1\); en el hiperbólico,
mediante \(\cosh^2t-\sinh^2t=1\); en el parabólico,
\((I+tN)(I-tN)=I\), pues \(N^2=0\). La misma ley de conservación
admite tres expresiones cuadráticas.

La realización elíptica cierra una vuelta; la nilpotente conserva un
desplazamiento lineal; la hiperbólica separa expansión y contracción con
determinante unitario. En los parámetros regionales ya generados,
\[
 \exp(\pi J_{\mathrm e})=-I,\qquad
 \exp(J_{\mathrm p})=I+J_{\mathrm p},\qquad
 \exp(2\log\varphi\,J_{\mathrm h})=F_{\rm av}^{\,2}.
\]
La evaluación \(\exp(I)=eI\) reconoce la unidad de la ley exponencial
común. El número \(e\) no se asigna exclusivamente al régimen parabólico.

El transporte de una realización local por un isomorfismo \(U\) conserva
estas relaciones: si \(J'=UJU^{-1}\), la serie exponencial da
\(U\exp(tJ)=\exp(tJ')U\). La identidad transporta el régimen y su
norma. Una transición entre tipos cuadráticos diferentes pertenece, en
cambio, a la selección de hojas y regímenes del TPK; no es una conjugación
que cambie el valor de \(J^2\). El producto temporal conserva por ello el
orden de las transiciones y la información de sus extremos.

Esta articulación distingue la clausura de una representación local y la
holonomía de una historia. Puede cumplirse \(\exp(2\pi J_{\mathrm e})=I\)
al mismo tiempo que \(S|r,m\rangle=|r,m+1\rangle\). La primera igualdad
restaura la orientación del régimen; la segunda registra otra vuelta. La
estructura discreta conserva ambas: el retorno local y la procedencia
acumulada del estado que lo realiza.
